% ===== §24 =====
\section{Fragility and Reversibility}
\label{p24:sec:fragility}

\noindent
The fifth pillar names the price of the other four. Everything that makes intimate culture what it is---the vacancy of the big Other that lets the couple author their own world, the interweaving of drive that gives their culture its incommunicable charge, the accelerated sedimentation that brings its emergence within experience, the intensity of the three registers met at close range---every one of these is also a condition of a fragility that no macro-culture faces. Intimate culture can collapse entire. It can pass, and often does, from a living world to a field of ruins, and the passage is not a gradual erosion but a structural reversal that the preceding pillars fully explain. This pillar describes that reversal and its mechanism, and in doing so completes the phenomenology of the dyad.

\subsection*{The asymmetry with macro-culture}

A macro-culture does not collapse when individuals leave it. Its forms are guaranteed by a big Other that stands above any individual---the tradition, the institution, the community of speakers---and the withdrawal or death of any member, or of many, leaves the culture's meanings in place, held by the super-individual structure that Simmel saw the dyad to lack. A language does not lose a word when a speaker dies; a rite does not cease to mean when a participant departs. The intimate dyad is, once again, different in kind, and the difference is the one the first pillar established. Because the couple's culture is guaranteed by no big Other above the two, but only by the two themselves, its entire existence is staked on the continued participation of both. The withdrawal of either member is not, as in macro-culture, the loss of one bearer among many; it is the removal of one of the only two guarantors there are, and since it takes both to hold the private meanings in place, the withdrawal of one does not halve the culture but ends it. This is Simmel's veto raised to the level of culture: as either member can end the dyad, so either member can end the culture the dyad generated, and the culture has no existence that survives the relation whose culture it was.

\subsection*{The mechanism of collapse: from meaning back to the Real}

The mechanism of the collapse can be stated precisely, in the terms the pillars have built, and it is more than the mere cessation of a shared activity. Recall that in the intimate dyad the place of the big Other, the authority that guarantees meaning, is occupied by the partner: each is, for the other, the guarantor of the private symbolic order they have built together. When one partner withdraws---through departure, betrayal, or death---the private symbolic order does not simply lose a user; it loses its guarantor, the very authority that held its signs in the position of signs. And a signifier that has lost the authority guaranteeing its meaning does not become a neutral, disused sign; it reverts. The private word, no longer guaranteed by the one who authorized it, falls from meaning back toward the Real: it becomes, first, a mere sound, emptied of the charge the interweaving gave it; and then, worse, a sound that signifies precisely the absence of the one who is gone, a marker of the void where the guarantor stood. The idioms that were the very substance of a shared world become, in the collapse, the sharpest instruments of grief, because each of them is now a private sign whose meaning was the bond and whose only remaining reference is the bond's destruction. This is why the ruins of an intimate culture are not neutral debris but actively painful: the private language does not fade into meaninglessness but is transformed, by the withdrawal of the partner who guaranteed it, into a language all of whose words now say the same thing, which is that the one who made them mean is gone. The collapse is a decoherence in the precise sense the framework gave the term: the symbolic forms remain, as sounds and gestures, while the real coherence that charged them is withdrawn, and what remains is the shell of a culture haunted by the bond it can no longer carry.

\subsection*{Exit costs and the asymmetry of collapse}

The collapse need not be symmetrical, and here the pillar connects to an analysis the series has developed elsewhere. The withdrawal that collapses the culture is not always mutual, and the two partners are not always equally exposed to the ruin. Where the exit costs are asymmetric---where one partner can leave at less cost than the other, whether because the interweaving was less deep on one side, or because the material and social conditions of the two differ, or because one has alternatives the other lacks---the collapse falls harder on the one for whom the exit cost was higher, who is left holding a private culture whose guarantor has departed at less cost to himself \parencite{paperxvii}. The one who leaves at low cost may re-enter the wider symbolic order, where public meanings are guaranteed by a big Other no single defection can unsettle; the one who leaves at high cost, or who does not leave but is left, remains among the ruins of a private symbolic order whose meanings have reverted to the Real. This asymmetry of exposure to collapse is the dyadic-cultural form of the asymmetry of power the series has analyzed under the heading of exit costs, and it shows that the fragility of intimate culture is not only a structural fact about dyads in general but a site at which the power asymmetries between two particular partners are registered and can be exploited.

\subsection*{Fragility as the price of endogeneity}

The pillar closes by returning fragility to its source, for it is essential that the fragility not be read as an unfortunate defect that a better-built intimate culture might avoid. The fragility is the price of the endogeneity, and the two are inseparable. A culture that a couple authors themselves, in the vacancy of any external authority, is by that very fact a culture with no external authority to preserve it when they falter; a culture whose charge is an interweaving of drive proper to two people is by that very fact a culture that cannot survive the unweaving of that pair; a culture sedimented fast, by two, over a rememberable time, is by that very fact a culture without the deep foundations that would let it outlast its makers. The freedom to make a world of one's own with another and the exposure to losing that world entire are the same condition seen from two sides, exactly as the possibility and the fragility were one in the first pillar. There is no intimate culture rich enough to be worth having that is not, by the very features that make it rich, exposed to collapse; the endogenous world and its fragility are bought together or not at all. This completes the phenomenology of the dyad's mechanisms, and the phenomenology turns now from the mechanisms to the interpretation of what they generate---to the hermeneutics of intimate culture, and to the negotiation, at its edge, with the outer symbolic order that the final part will show to be the site of its capture.
